% Part: computability
% Chapter: computability-theory
% Section: complete-ce-sets

\documentclass[../../../include/open-logic-section]{subfiles}

\begin{document}

\olfileid{cmp}{thy}{cce}
\olsection{المجموعات التامة القابلة للتعداد حسابيًا}

\begin{defn}
نسمي ${\OLId{latin-looped}{A}}$ \emph{مجموعة تامة !!{computably enumerable}}
بالنسبة إلى الاختزال متعددًا إلى واحد، إذا اجتمع الشرطان:
\begin{enumerate}
\item كانت ${\OLId{latin-looped}{A}}$ قابلة للتعداد حسابيًا، و
\item كان ${\OLId{latin-looped}{B}}\leq_{\OLId{latin-ordinary}{m}} {\OLId{latin-looped}{A}}$ لكل مجموعة أخرى ${\OLId{latin-looped}{B}}$ قابلة للتعداد حسابيًا.
\end{enumerate}
\end{defn}

فهذه المجموعات التامة هي «أصعب» ما يمكن من
المجموعات القابلة للتعداد حسابيًا؛ إذ بمعرفة أجوبتها نجيب عن الأسئلة
المتعلقة بـ\emph{أي} مجموعة قابلة للتعداد حسابيًا.

\begin{thm}
${\OLId{latin-looped}{K}}$ و${\OLId{latin-looped}{K}}_{\OLMathNumeral{0}}$ و${\OLId{latin-looped}{K}}_{\OLMathNumeral{1}}$ كلها مجموعات تامة قابلة للتعداد حسابيًا.
\end{thm}

\begin{proof}
أما تمام ${\OLId{latin-looped}{K}}_{\OLMathNumeral{0}}$، فلنأخذ ${\OLId{latin-looped}{B}}$ قابلة للتعداد حسابيًا كيفما اتفق؛
فيكون لها فهرس ${\OLId{latin-ordinary}{e}}$ يحقق
\[
{\OLId{latin-looped}{B}} = {\OLId{latin-looped}{W}}_{\OLId{latin-ordinary}{e}} = \Setabs{{\OLId{latin-ordinary}{x}}}{\cfind{{\OLId{latin-ordinary}{e}}}({\OLId{latin-ordinary}{x}}) \fdefined}.
\]
ونختار ${\OLId{latin-ordinary}{f}}$ بحيث ${\OLId{latin-ordinary}{f}}({\OLId{latin-ordinary}{x}})=\tuple{{\OLId{latin-ordinary}{e}},{\OLId{latin-ordinary}{x}}}$. فلكل عدد طبيعي ${\OLId{latin-ordinary}{x}}$ يتحقق
${\OLId{latin-ordinary}{x}}\in {\OLId{latin-looped}{B}}$ إذا وفقط إذا كان ${\OLId{latin-ordinary}{f}}({\OLId{latin-ordinary}{x}})\in {\OLId{latin-looped}{K}}_{\OLMathNumeral{0}}$. فتكون ${\OLId{latin-ordinary}{f}}$
اختزالًا من ${\OLId{latin-looped}{B}}$ إلى~${\OLId{latin-looped}{K}}_{\OLMathNumeral{0}}$.

وأما ${\OLId{latin-looped}{K}}_{\OLMathNumeral{1}}$ فقد اختزلنا ${\OLId{latin-looped}{K}}_{\OLMathNumeral{0}}$ إليها في برهان
\olref[k1]{prop:k1}؛ وبالتعدي الوارد في \olref[ppr]{prop:trans-red}،
تختزل كل مجموعة قابلة للتعداد حسابيًا إلى~${\OLId{latin-looped}{K}}_{\OLMathNumeral{1}}$ أيضًا.

وبمثل هذا الطريق نختزل ${\OLId{latin-looped}{K}}_{\OLMathNumeral{0}}$ إلى ${\OLId{latin-looped}{K}}$.
\end{proof}

\begin{prob}
بيّن اختزال ${\OLId{latin-looped}{K}}_{\OLMathNumeral{0}}$ إلى ${\OLId{latin-looped}{K}}$.
\end{prob}

\begin{digress}
فالأمثلة التي مرت بنا إلى الآن من المجموعات القابلة للتعداد حسابيًا
كلها إما قابلة للحساب وإما تامة. وهذا موضع سؤال:{}
هل بين الطرفين مجموعات قابلة للتعداد حسابيًا، لا هي قابلة للحساب ولا
هي تامة؟ نعم؛ وينسب النص الأصلي إثبات ذلك إلى فريدبرغ وموشنيك، إذ توصلا إليه مستقلين في
منتصف خمسينيات القرن العشرين.
\end{digress}

\end{document}
